\section{Introduction}
\label{sec:intro}
Recent developments in natural language processing (NLP) have shifted the paradigm from monolithic language models to modular, multi-agent frameworks \citep{jenkins2024dsr}. In these environments, distinct agents are specialized in specific sub-domains of linguistic reasoning, such as syntax, semantics, and pragmatics. However, coordinates for organizing these agents efficiently remain a persistent bottleneck. Static routing policies often lead to under-utilization or cognitive overload of specialized agents.

To address this challenge, we introduce the Deep Synaptic Routing (DSR) framework. DSR dynamically directs context representations to specific linguistic agents based on a gating mechanism that monitors the input sequence. Unlike traditional Mixture-of-Experts (MoE) methods, which assign tokens individually \citep{vance2023neural}, DSR operates at the level of semantic units, maintaining semantic coherence across sentences.

The primary contributions of this paper are:
\begin{enumerate}
  \item We formulate a dynamic routing objective function that minimizes classification error and routing latency simultaneously.
  \item We show that our router constructs a latent space where agents form functional hierarchies.
  \item We provide empirical evaluations demonstrating the robustness of DSR on multi-agent translation and reasoning benchmarks.
\end{enumerate}
